\section{Discussion and Future Work}
\label{sec:discussion}

To the best of our knowledge, Continual-ARC is the first benchmark to compare harness-based and weight-based learners under one score, in a domain where both can memorize and generalize. We have shown that most harnesses as shipped did not learn measurably in this domain. How well they learn in other domains remains open. Simple expert guidance greatly improved how much they memorized (Figure~\ref{fig:plane-luna}), which shows that the default harness prompts were not sufficient to inform the models what to store, though more powerful models may fare better.

Training the weights proved important in this domain, as DAgger learned from a short stream of tasks to solve new ones where in-context learning through the harness did not. Continual-ARC targets the continual learning that one person or a small team meets over a short horizon, from 200 instances rather than a post-training corpus. We are therefore most excited about weight-based algorithms that learn from sparse data, model-based reinforcement learning above all \citep{kaiser2020model}. The DAgger algorithm we used is a relatively simple form of weight-based learning, as it needs a stronger teacher whose reasoning tokens are available. Teacher-free reinforcement learning from the sparse reward \citep{shao2024deepseekmath}, on-policy self-distillation \citep{shenfeld2026sdft,hubotter2026sdpo} and on-policy distillation \citep{agarwal2024gkd} are natural next steps (Appendix~\ref{app:other-weights} reports our first attempts). Trust-region updates \citep{schulman2015trpo} would also keep such a learner from degrading what its users receive, which is essential in real deployments.

One weakness of this study is that most of its runs used relatively small models on ARC-AGI-1, the easiest of our three tracks. We leave it to future work to test stronger models on the ARC-AGI-2 and ARC-AGI-3 tracks (Appendix~\ref{app:tracks}). This will let us investigate how much guidance frontier models need in order to store what matters in a harness, and how efficiently they learn it in their weights, both for memorization and generalization. In the long run, we expect the harness to become less the system itself and more the infrastructure that trains the inner model, continually distilling content into the weights and staying thin rather than accumulating an ever-growing list of files.
